language at infinity: the asymptotic anatomy of linguistic systems

the useful question is not merely, “what does a very long sentence look like?” it is: when some measure of linguistic size tends toward infinity, which properties stabilize, which diverge, which undergo phase transitions, and which details disappear into a smaller number of universality classes?

this is analogous to studying infinitely wide neural networks, thermodynamic limits, or high-dimensional probability—but with one major complication. there is no single scalar called “the size of language.” at least six different quantities can independently tend to infinity:

[
\text{utterance length }n,\quad
\text{grammar size }|g|,\quad
\text{corpus size }n_c,\quad
\text{semantic graph size }|s|,\quad
\text{dialogue horizon }t,\quad
\text{speaker population }p.
]

these limits do not commute. taking sentence length to infinity while keeping the grammar fixed is not equivalent to enlarging the grammar while keeping sentences short. taking the corpus to infinity reveals the distribution induced by a language community, while taking the semantic world to infinity concerns how finite expressions select states from an unbounded reality. taking the dialogue horizon to infinity produces a dynamical system, not merely a larger syntax tree.

therefore, before asking about “the infinite language tree,” we need to specify:

[
\text{object}
+\text{sampling law}
+\text{conditioning}
+\text{rescaling}
+\text{equivalence relation}.
]

without those five elements, the limit is underdetermined.

there is another correction worth making immediately. an lr parser is a recognition procedure. its state stack and transition graph are properties of a grammar presentation and parsing algorithm, not intrinsic properties of the language itself. two grammars can generate the same strings while producing very different lr automata, derivation trees, ambiguity patterns, and stack trajectories. so an infinite lr tree would primarily characterize a parser–grammar pair.

for the deeper linguistic question, better primitive objects are:

[
x_{1:n} \quad \text{the token sequence},
]

[
t_n \quad \text{a constituency or dependency tree},
]

[
s_n \quad \text{a sentence-level semantic graph},
]

[
d_n \quad \text{a discourse graph},
]

[
w_t \quad \text{a latent world or common-ground state}.
]

the same utterance should normally be analyzed through all five. language is not secretly “a tree.” it is a family of mutually constraining projections of a generative, semantic, social, and causal process.

## the first limit: the derivation tree

suppose a grammatical constituent of type (i) produces a collection of child constituents of types (j). abstractly, this resembles a multi-type branching process. let

[
m_{ij}
======

\mathbb{e}[\text{number of type-}j\text{ children generated by a type-}i\text{ node}],
]

and let (m) be the resulting reproduction matrix. the spectral radius

[
\rho(m)
]

divides the unconditioned branching process into three broad regimes:

[
\rho(m)<1
\quad \text{subcritical},
]

[
\rho(m)=1
\quad \text{critical},
]

[
\rho(m)>1
\quad \text{supercritical}.
]

subcritical trees tend to terminate rapidly. supercritical trees have a positive probability of growing indefinitely. critical trees sit at the boundary: they terminate almost surely under ordinary assumptions, but can generate arbitrarily large structures with heavy-tailed size distributions.

natural utterances are finite, so a useful construction is not simply “run the branching grammar forever.” instead, sample a derivation tree conditioned on having exactly (n) terminal or total nodes, then ask what happens as (n\to\infty).

this is where random-tree theory gives us a striking distinction between global and local limits.

for broad classes of critical, finite-variance conditioned galton–watson trees, total height and width are typically on the order of (\sqrt n). after graph distances are rescaled by (n^{-1/2}), the global tree can converge toward a continuum random tree encoded by a brownian excursion. from far enough away, many microscopically different branching rules therefore collapse into the same macroscopic geometry.

but if we keep distances unscaled and inspect a fixed-radius neighborhood near the root, the limit can instead be an infinite tree with a distinguished endless spine and finite side branches: a kesten-type tree. this means that the same sequence of finite trees can have a compact continuum global limit and an infinite-spine local limit. neither is “the” limit; they answer different questions.

a third regime is condensation. under certain non-generic or heavy-tailed conditions, a single node acquires a macroscopic fraction of all children. its limiting local structure contains a node of effectively infinite degree rather than an infinite spine.

these three shapes already suggest useful hypotheses for language:

a spine-dominated structure corresponds to repeated continuation through one persistent recursive thread, with bounded elaborations attached along it.

a continuum-critical structure corresponds to balanced recursive elaboration spread across many scales.

a condensation structure corresponds to one discourse topic, conjunction, list, event, or governing predicate accumulating an extensive number of dependents.

these are not claims that ordinary syntax literally follows a galton–watson law. they are candidate universality classes against which empirical syntax trees can be compared.

## the asymptotic phenotype of a syntax tree

the first property is the language’s combinatorial growth rate. let

[
a_n
===

|{x\in l:|x|=n}|,
]

the number of distinct length-(n) strings admitted by a formal language. then

[
h_{\mathrm{top}}
================

\limsup_{n\to\infty}\frac{1}{n}\log a_n
]

is a topological entropy: the exponential rate at which the language’s support expands.

two languages may have the same (h_{\mathrm{top}}) but very different subexponential structure. a more precise asymptotic form might be

[
a_n
\sim
c\mu^n n^{-\alpha}.
]

here (\mu=e^{h_{\mathrm{top}}}) controls the exponential growth, while (\alpha) often reveals information about recursive constraints and singularities in the grammar’s generating function. this is already richer than vocabulary size or perplexity. it distinguishes languages whose option spaces expand at the same exponential rate but whose long strings occupy very different geometric fractions of that space.

the second property is the height exponent:

[
h_n\asymp n^{\beta_h}.
]

a chain has (\beta_h=1). a balanced binary tree has logarithmic height, which is effectively (\beta_h=0) with a logarithmic correction. many critical random trees have (\beta_h=\tfrac12). measuring (\beta_h) would tell us whether additional linguistic material is incorporated primarily through sequential embedding, balanced subdivision, or broad attachment.

the third property is the width profile. define

[
w_n(k)
======

|{v\in t_n:\operatorname{depth}(v)=k}|.
]

the entire function (w_n(k)), after horizontal and vertical rescaling, may converge to a stable limit shape. two registers could have identical average depth but very different profiles: one may expand abruptly near the root and then taper; another may remain narrow for a long time and branch only near the leaves.

the fourth property is imbalance. for a binary node with subtree sizes (n_l) and (n_r), one elementary measure is

[
b(v)
====

\frac{|n_l-n_r|}{n_l+n_r}.
]

averaging this over nodes, or examining its depth-dependent distribution, distinguishes right- or left-branching chains from recursively balanced explanation. imbalance matters computationally because it changes stack depth, incremental interpretability, and the distance between constituents that must remain jointly active.

the fifth property is branching-direction persistence. even with the same degree distribution and height, a tree may repeatedly branch in the same direction or alternate directions. this captures the difference between persistent center embedding, repeated right adjunction, and more symmetric recursive decomposition.

the sixth is the degree tail:

[
p(\deg v\ge k).
]

a thin tail implies a hard or effective upper bound on local fan-out. a heavy tail permits occasional large coordination, enumeration, or modifier bundles. a condensation parameter could quantify whether the largest degree remains negligible or captures an extensive fraction of the tree:

[
c_n
===

\frac{\max_v \deg(v)}{|t_n|}.
]

if (c_n\to0), there is no macroscopic hub. if (c_n\to c>0), a finite number of constituents dominate global structure.

the seventh is fringe convergence. choose a random node and examine the finite subtree beneath it. as (n) grows, does the distribution of these local patterns stabilize? this reveals whether large texts are composed from a stationary inventory of local constructions or whether local grammar itself changes systematically with document scale.

the eighth is motif frequency. one can count relative clauses, coordination motifs, argument frames, modifier chains, repair fragments, quotation embeddings, or any other typed subtree. the limit

[
f_m
===

\lim_{n\to\infty}
\frac{#\text{ occurrences of motif }m}{n}
]

is a grammatical density rather than a mere corpus count. its finite-size convergence rate is also informative: stable function-word constructions may converge quickly, while rare rhetorical structures require document-level scale.

the ninth is the tree’s strahler number or register complexity. this asks how many recursively comparable branches must coexist before the tree can be reduced. it distinguishes a very long chain from a shorter but genuinely parallel hierarchy. both can have large depth, but only the latter requires maintaining multiple complex subproblems simultaneously.

the tenth is crossing complexity. dependency trees are often drawn above a token sequence, giving each edge a linear span. one can measure crossing count, maximum crossing degree, gap degree, or the genus of a derived embedding. this characterizes how severely hierarchical adjacency diverges from sequential adjacency.

the eleventh is dependency-length scaling:

[
d_n
===

\frac{1}{|e_n|}
\sum_{(i,j)\in e_n}|i-j|.
]

the scaling of (d_n) with sentence length separates languages that keep related items locally adjacent from languages that tolerate increasingly nonlocal dependencies. this is not only a syntactic property; it constrains working memory and transmission robustness.

the twelfth is ambiguity growth. let (a(x)) be the number of valid derivations of string (x). then we can ask whether

[
\mathbb{e}[\log a(x)\mid |x|=n]
]

is bounded, logarithmic, sublinear, or linear in (n). linear growth means the number of parses is exponentially large. two languages with nearly identical string distributions could have radically different latent ambiguity.

but raw parse count treats trivial and meaningful alternatives equally. a stronger measure is posterior parse entropy:

[
h(t\mid x)
==========

-\sum_t p(t\mid x)\log p(t\mid x).
]

this measures how much structural uncertainty remains after observing the utterance.

## from trees to graphs

trees enforce a unique parent for every non-root node. semantics often does not.

the same entity can be an argument of several events. a person can be the agent of one action, the patient of another, the holder of a belief, and the referent of a pronoun. shared arguments produce reentrancy: multiple incoming paths to one semantic node. this is one reason systems such as abstract meaning representation use directed graphs rather than pure trees.

discourse adds more violations of tree structure:

coreference creates long-range identity edges.

causal relations create chains and converging causes.

temporal relations can create partial orders.

belief and quotation introduce nested contexts.

definitions create cycles among concepts.

presuppositions connect an utterance to prior commitments.

analogy connects otherwise distant semantic regions.

contradictions create typed incompatibility edges rather than compositional parent–child edges.

therefore, an infinite semantic limit should be characterized by graph observables.

the first is volume growth:

[
v(r)
====

|{u:d(u,v)\le r}|.
]

if (v(r)) grows polynomially, one can define an effective dimension (d) through (v(r)\sim r^d). if it grows exponentially, the graph is more tree-like or negatively curved. volume growth tells us how many meanings become accessible after a bounded number of inferential steps.

the second is average path length and diameter. a graph can be extremely sparse while preserving short routes between almost every pair of concepts. empirical association networks, wordnet-derived networks, and thesaurus networks have been reported to show sparse connectivity, high clustering, short paths, and hub-dominated degree distributions. these properties depend strongly on how edges are constructed, but they demonstrate that the large-scale semantic object need not resemble either a pure hierarchy or an unstructured vector cloud.

the third is clustering:

[
c(v)
====

\frac{\text{edges among neighbors of }v}
{\text{possible edges among neighbors of }v}.
]

high clustering means local semantic neighborhoods form coherent subdomains. low global path length combined with high clustering is the conventional small-world pattern.

the fourth is modularity. technical disciplines should induce dense subgraphs connected by narrower interfaces. “tensor,” “gradient,” and “backpropagation” lie in a machine-learning module; “standing,” “jurisdiction,” and “remedy” lie in a legal module. interdisciplinary discourse can be studied through the number, strength, and semantic stability of edges crossing such modules.

the fifth is conductance. if a semantic region has low conductance, random or associative motion tends to remain trapped there. low-conductance modules correspond to conceptual domains that require rare bridge concepts to exit. an expert may know precisely which bridge nodes connect two fields; a novice may possess many local terms while lacking these inter-domain edges.

the sixth is hub structure. generic relations such as cause, part, type, use, change, and possess may connect many domains. however, high degree can mean several different things: polysemy, abstraction, grammatical utility, frequency, or genuine causal centrality. a useful representation must type the edges before interpreting hubs.

the seventh is treewidth. a tree has treewidth one. reentrant, mutually constraining semantic structures have greater treewidth. intuitively, treewidth measures how much of the graph must be considered jointly to solve inference exactly through variable elimination. technical language describing tightly coupled systems may have moderate syntax-tree complexity but high semantic treewidth.

the eighth is cycle rank. ontologies often attempt to remain acyclic under “is-a” relations, but realistic knowledge graphs include feedback, mutual definition, regulation loops, and circular dependency. cycle statistics distinguish merely hierarchical taxonomies from dynamical models.

the ninth is hyperbolicity. negatively curved spaces efficiently represent exponentially branching hierarchies. a semantic graph’s hyperbolicity parameter can quantify whether triangles are thin and routes rapidly diverge as they move away from a common ancestor. this may distinguish taxonomic technical language from associative or narrative language, where lateral shortcuts are more common.

the tenth is spectral structure. eigenvalues and eigenvectors of a graph laplacian reveal bottlenecks, diffusion rates, communities, and multi-scale organization. rather than asking only whether two words have similar embeddings, one could compare the spectrum of entire induced semantic graphs.

the eleventh is spectral dimension. consider a random walk beginning at a semantic node. if the return probability scales as

[
p_{\mathrm{return}}(t)\sim t^{-d_s/2},
]

then (d_s) acts as a diffusion dimension. it describes how quickly semantic activation disperses, which need not equal the graph’s volume-growth dimension.

the twelfth is symmetry. an automorphism is a permutation of graph nodes that preserves all typed relations. a graph with many automorphisms contains interchangeable structural roles. language learning should ideally respect these symmetries: renaming entities in a structurally identical world should not alter the inferred relation pattern.

this leads to the most important semantic issue: equivalence.

## the quotient problem

a raw language graph contains distinctions that may be representational rather than semantic.

“alice gave bob the key.”

“bob received the key from alice.”

these expressions differ syntactically but may describe nearly the same event. to study semantic structure, we need an equivalence relation

[
x\sim y
]

meaning that (x) and (y) are identical for the level of analysis we care about.

but “semantic equivalence” is never completely context-free. the two sentences may differ in information structure, discourse focus, presupposition, causal framing, or expected continuation. therefore there is not one universal quotient. there is a hierarchy:

[
\text{truth-conditional equivalence},
]

[
\text{event-structural equivalence},
]

[
\text{causal equivalence},
]

[
\text{pragmatic equivalence},
]

[
\text{action-policy equivalence}.
]

the limiting semantic graph depends on which quotient is chosen.

this creates a gauge freedom. two internal representations may use different latent node names, coordinate systems, parse decompositions, or intermediate propositions while inducing the same observable predictions and actions. the representation is identifiable only up to transformations that preserve observables.

therefore, merely probing whether a neural network contains a node corresponding to “agent” or “object” is weaker than testing whether its representation is equivariant under entity permutation, paraphrase, syntactic alternation, and causal intervention.

a genuine theory of language at infinity should separate:

[
\text{intrinsic invariants}
]

from

[
\text{artifacts of the selected parser or coordinate system}.
]

## information at infinity

the token-level entropy rate is

[
h
=

\lim_{n\to\infty}
h(x_n\mid x_{1:n-1}).
]

it measures the irreducible uncertainty per additional token when arbitrarily long context is available. neural language-model cross-entropy is an upper bound or approximation to this quantity, contaminated by finite model capacity, finite data, optimization error, tokenization, and context truncation.

entropy rate alone misses almost everything about hierarchy. two processes can have the same entropy rate but distribute information across time in completely different ways.

we therefore need the excess entropy or predictive information between large neighboring blocks:

[
e(n)
====

i(x_{1:n};x_{n+1:2n}).
]

if (e(n)) converges to a constant, distant past contributes only bounded total predictive structure. if it grows logarithmically or as a power law, progressively larger scales reveal additional shared information. work on natural-language excess entropy has explored how power-law information growth can relate vocabulary growth and texts that repeatedly describe a complex latent object.

another observable is mutual information at separation (d):

[
i(d)
====

i(x_t;x_{t+d}).
]

possible asymptotic regimes include

[
i(d)\sim e^{-d/\xi},
]

[
i(d)\sim d^{-\eta},
]

or slower, multi-regime decays.

the exponential regime has a characteristic correlation length (\xi). beyond several multiples of (\xi), direct dependence is negligible.

a power law has no single characteristic scale. structure persists across many distances, which is one mathematical sense in which a system can be scale-free.

lin and tegmark showed that probabilistic regular grammars have exponentially decaying symbol mutual information, while probabilistic context-free constructions can produce power-law decay through recursive hierarchy.

however, that does not establish that observed power laws in human text are caused by context-free syntax. recent corpus analysis of english and japanese parse trees found that key assumptions needed for the simple pcfg explanation fail empirically, including substantial violations of context-free independence at large structural distances. the general lesson is crucial: a power law is a phenotype, not a unique mechanistic diagnosis.

other important information-theoretic observables include:

the surprisal distribution

[
s_t=-\log p(x_t\mid x_{<t}),
]

not merely its mean.

surprisal variance and tail behavior, which distinguish steady information flow from intermittent bursts.

multi-information among more than two separated regions.

conditional mutual information after removing topic, author, or document identity.

directed information, measuring whether earlier structure predicts later structure asymmetrically.

semantic entropy, where alternative surface realizations of the same meaning are collapsed.

ambiguity entropy, measuring uncertainty over meanings for a fixed string.

revision entropy, measuring how often later context changes the interpretation of an earlier span.

compression depth, measuring how much intermediate latent structure is required to achieve near-optimal prediction.

a strong asymptotic description should also separate support complexity from probability complexity. (h_{\mathrm{top}}) measures how many strings are possible. (h) measures how probability is distributed over those possibilities. a language can admit an enormous support while concentrating almost all probability on a tiny conventional submanifold.

## the latent-world limit

a sentence is not merely generated from a grammar. it updates a posterior over possible worlds:

[
p(w\mid x_{1:t},c).
]

as tokens arrive, some worlds are eliminated, others split into finer alternatives, and some that appeared distinct become observationally equivalent.

this suggests a semantic branching process that is almost dual to the syntax tree. the syntax tree expands as the utterance is generated; the posterior world set contracts as evidence accumulates.

define the effective number of plausible worlds by posterior perplexity:

[
n_{\mathrm{world}}(t)
=====================

\exp h(w\mid x_{1:t},c).
]

the ambiguity-collapse curve

[
a(t)
====

h(w\mid x_{1:t},c)
]

tells us how quickly an utterance resolves uncertainty.

technical terminology may cause sharp entropy drops: once a precise term is recognized within the correct domain, many interpretations disappear simultaneously.

poetic language may deliberately maintain a high-entropy posterior for longer.

legal drafting attempts to eliminate some dimensions of uncertainty while occasionally leaving strategic uncertainty elsewhere.

casual dialogue often has low surface explicitness but can still produce rapid posterior contraction because the participants share visual, social, and historical context.

we can also define an ambiguity survival function. let two world hypotheses (w_i,w_j) remain meaningfully distinct until token (t). then

[
q(d)
====

p(\text{a distinction survives for at least }d\text{ tokens})
]

measures the lifetime distribution of unresolved interpretations.

another property is entity persistence. construct a node for each discourse entity and track the distance between successive mentions or inferred activations. an entity half-life measures how quickly it leaves the active discourse state.

causal depth is the longest relevant chain of interventions required to justify a proposition:

[
u_0\rightarrow u_1\rightarrow\cdots\rightarrow u_k.
]

this can be much greater than syntactic depth. “the chip failed because contamination altered the oxide growth, which shifted threshold voltage” may have fairly ordinary syntax while encoding a multi-level causal path.

counterfactual dimension measures the number of independently variable interventions needed to distinguish the worlds compatible with an utterance.

epistemic depth measures nested belief structure:

[
\text{alice believes that bob suspects that carol knows }p.
]

modal breadth measures the number of possible, necessary, permitted, intended, or expected branches represented simultaneously.

inferential diameter is the maximum or typical number of rule applications needed to connect claims that the discourse treats as mutually relevant.

contradiction topology records which propositions cannot coexist, under what assumptions, and whether contradictions remain local or propagate globally.

these properties are nearly invisible to ordinary next-token perplexity.

## dialogue at infinity

when dialogue length tends to infinity, the correct object is a coupled dynamical system.

let

[
c_t
]

represent common ground at turn (t), including discourse entities, commitments, unresolved questions, goals, social state, and beliefs about each participant’s beliefs. each utterance applies an update

[
c_{t+1}
=======

f(c_t,u_t,a_t,e_t),
]

where (u_t) is the linguistic signal, (a_t) is the accompanying action, and (e_t) is the external environment.

spoken dialogue is routinely elliptical, anaphoric, incrementally interpreted, and repaired across participants. semantic construction can cross turn boundaries, and systems must revise interpretations before an utterance has even formally ended.

an infinite-horizon dialogue can exhibit several dynamical regimes.

it may converge toward an absorbing state: a question is answered, a plan is accepted, or disagreement is resolved.

it may enter a limit cycle: the participants repeat the same argument or clarification pattern.

it may remain recurrent but nonperiodic, revisiting themes with altered local state.

it may drift continually as new facts and goals enter.

it may branch explosively, generating more unresolved commitments than it closes.

define the unresolved-question load (q_t). then the drift

[
\Delta q
========

\lim_{t\to\infty}
\frac{\mathbb{e}[q_t-q_0]}{t}
]

distinguishes convergent dialogues from conversations that accumulate semantic debt.

the repair rate measures how often participants retract, replace, or reinterpret prior structure.

the repair radius measures how far backward a repair reaches.

the turn-coupling coefficient measures how much information in one speaker’s utterance can only be interpreted using the immediately preceding contribution.

the exophoric load measures how much meaning is recovered from the physical or social environment rather than the linguistic transcript.

the common-ground compression ratio measures how many tokens can be omitted because participants already share a latent model.

the response branching factor measures how many pragmatically acceptable continuations remain after each turn.

the convergence rate measures how rapidly participant belief states become aligned—or at least mutually predictable.

these dynamics are central to distinguishing dialogue registers. a transcript can appear syntactically simple precisely because most of its structure exists in common ground rather than in the words.

## asymptotic fingerprints of linguistic registers

the interesting possibility is that language types are not best distinguished by a single complexity score. they may occupy different regions of a high-dimensional asymptotic phase space.

### expert-to-expert technical dialogue

expert language is highly compressed relative to a shared ontology.

within the correct domain, a technical term can have low semantic branching: it sharply selects one concept and its expected relation structure. but the term may sit at the end of a very deep prerequisite path. the utterance is locally precise and globally inaccessible.

expert dialogue should therefore often exhibit:

high presupposed-graph size.

high common-ground compression.

high terminology consistency.

low within-domain synonym freedom.

dense typed semantic relations.

high causal and inferential depth.

strong modularity.

rapid posterior contraction after key terminology.

long entity persistence.

relatively low explicit redundancy.

this does not imply universally deeper syntax. expert writing varies heavily by field, and technical material may relocate complexity into equations, diagrams, code, tables, or shared formalism. one corpus study found substantial differences in lexical and syntactic complexity across economics, it, linguistics, and political science; it texts were comparatively simple by textual metrics, with some omitted complexity residing in formulas, code, and symbols.

expertise is therefore better described as semantic compression backed by a large shared graph, not as “using longer words and sentences.”

### scientific exposition

scientific prose attempts to be self-contained enough for asynchronous inspection, but it still assumes disciplinary background.

likely properties include:

dense noun-phrase modification.

high reference and citation connectivity.

low tolerance for uncontrolled lexical drift.

explicit evidential and epistemic typing.

high method–result causal structure.

moderate redundancy around central claims.

long-range coreference across sections.

strong discourse segmentation into motivation, method, result, and qualification.

genre studies find that scientific and administrative writing make heavy use of dense noun modification to encode technical details and precision, while more conversational or literary genres rely more strongly on contextual inference and ellipsis.

scientific language may therefore have lower surface branching than literary text but greater semantic treewidth: many variables, assumptions, and measurements constrain one another simultaneously.

### pedagogical explanation

pedagogical language deliberately reverses some expert compression.

a good lesson expands edges that experts normally traverse implicitly. it supplies intermediate nodes, repeats a concept through multiple encodings, gives examples and counterexamples, and regularly reconnects local details to a global schema.

its asymptotic signature should include:

lower prerequisite depth per unexplained term.

higher semantic redundancy.

multiple paraphrase paths to the same concept.

lower ambiguity survival.

controlled growth of syntactic depth.

high cross-scale recurrence of central ideas.

frequent local summaries.

explicit bridge nodes between conceptual modules.

the optimal pedagogical graph is not simply smaller. it has greater navigational robustness: if one explanatory path fails, another reaches the same semantic region.

### vernacular conversation

ordinary conversation externalizes relatively little because participants share the room, recent history, social expectations, and perceptual context.

it tends toward:

short syntactic fragments.

high turn coupling.

high ellipsis and anaphora.

high exophoric load.

frequent repair.

rapid topic switching.

socially meaningful lexical variation.

low transcript-level self-containment.

high dependence on prosody, gesture, timing, and gaze.

its explicit syntax graph may be sparse, while its full multimodal semantic graph is dense. treating only the words as the language object makes vernacular conversation falsely appear impoverished.

### legal and administrative language

legal language is shaped by adversarial interpretation and downstream enforcement.

its characteristic properties may include:

high quantifier density.

high modal density.

large scope dependencies.

many exceptions and conditional branches.

low permissible paraphrase freedom.

high reference persistence.

explicit definition subgraphs.

low tolerance for contextual ellipsis.

high contradiction sensitivity.

potentially large semantic treewidth.

administrative and legal-style prose often uses dense nominal structures and noun modifiers to encode precision, though the resulting surface form can become difficult to process.

its complexity is not merely depth. a contract often resembles a constraint system whose clauses jointly narrow the admissible world set.

### literary and poetic language

literary language can optimize for controlled underdetermination rather than immediate posterior collapse.

its likely properties include:

high semantic ambiguity entropy.

slow ambiguity decay.

long-range motif recurrence.

cross-domain shortcuts.

unusual lexical neighborhoods.

high sensitivity to exact surface form.

multiple valid semantic quotients.

deliberate violation of ordinary surprisal profiles.

persistent unresolved entities or causal relations.

poetry may compress syntax through ellipsis while expanding the interpretation graph. this is nearly the opposite of technical jargon, which can expand the surface term while collapsing interpretation within a trained community.

### social-media vernacular and slang

this register combines short utterance horizons with rapid population-level evolution.

relevant limits include speaker population and time, not merely sentence length.

its phenotype may include:

high lexical birth and death rates.

strong community modularity.

high indexical and identity information.

short-lived semantic conventions.

rapid edge rewiring.

high dependence on external events.

extreme compression under shared memes.

large meaning shifts from tiny orthographic changes.

the important property is graph velocity:

[
v_g
===

\frac{\text{semantic edge additions, deletions, and relabelings}}
{\text{time}}.
]

a static dictionary cannot describe a language regime whose semantic topology changes on the same timescale as its use.

### mathematics and programming languages

formal languages intentionally collapse ambiguity at the cost of requiring more explicit structure.

they tend toward:

low parse entropy.

strong compositionality.

exact symbol identity.

large possible recursion depth.

strictly typed relations.

high sensitivity to local errors.

executable or mechanically checkable semantics.

but actual mathematical and programming practice remains hybrid. comments, naming, theorem motivation, api expectations, and problem selection remain natural-language and world-model phenomena.

formal code can be syntactically unambiguous while semantically underdetermined with respect to the programmer’s intention. therefore execution equivalence, denotational equivalence, and intention equivalence define different quotients.

## what current neural networks measure

current language models already provide imperfect proxies for some asymptotic properties.

cross-entropy approximates entropy rate.

performance as context length increases approximates a finite-memory curve.

syntactic probes estimate whether local grammatical relations are linearly recoverable.

semantic parsing evaluates reconstruction of selected graph schemas.

embedding distances provide a weak view of semantic geometry.

activation patching and causal interventions test whether internal states contribute to outputs.

scaling laws describe how loss changes with parameters, data, and compute.

but these proxies are mostly point estimates. they do not characterize the induced linguistic object as (n\to\infty).

attention distance is not a dependency-length measure unless attention has been causally shown to implement that dependency.

a linear probe finding syntactic labels does not establish that the model computes through a syntax tree.

low perplexity does not imply correct ambiguity management.

good paraphrase similarity does not imply causal or pragmatic equivalence.

successful next-token prediction does not reveal whether the model maintains one world, a posterior over worlds, or a lossy superposition that only works on familiar continuations.

## properties for which we largely lack neural proxies

the first missing proxy is latent limit shape. as input length increases, extract a model’s effective computational or relational graph, rescale it, and test whether its depth, width, spectrum, and motif distributions converge.

the second is local weak convergence. choose a random token or concept node and inspect a fixed-radius neighborhood in the model’s induced graph. does its distribution stabilize as documents grow?

the third is ambiguity survival. require the model to preserve several incompatible but viable interpretations and measure when each branch is correctly eliminated. ordinary language models are rewarded for assigning token probabilities, not for maintaining an explicit, calibrated population of latent worlds.

the fourth is semantic quotient invariance. generate meaning-preserving graph transformations—active/passive alternation, entity renaming, argument permutation with corresponding role changes, equivalent logical rewrites—and test whether the internal representation transforms predictably.

the fifth is intervention equivariance. modify the latent causal world while preserving irrelevant surface statistics. the model’s answers should change only along the causal descendants of the intervention.

the sixth is topological support entropy. rather than measuring probability on observed text, estimate how the number of grammatical, coherent, or task-valid continuations grows with length.

the seventh is induced treewidth. determine how many interacting semantic variables the model can maintain without factorizing them incorrectly.

the eighth is graph-surgery generalization. train on local semantic motifs, then combine them into larger graphs with familiar components but novel global topology. this tests whether the model learns reusable relations or memorizes common whole structures.

the ninth is asymptotic repair radius. insert contradictions or corrections at increasing distances and measure whether the model properly revises all downstream implications.

the tenth is world-state persistence. introduce entities and latent attributes, allow them to disappear from the surface for increasingly long intervals, and test whether identity and state remain coherent.

the eleventh is cross-scale information matching. compare the mutual-information and excess-entropy curves of model-generated text against human corpora, not merely unigram frequencies or short-context likelihoods.

the twelfth is register phase classification. estimate a full vector of asymptotic properties and test whether expert dialogue, teaching, bureaucracy, casual speech, literature, and code form distinct clusters or continuous manifolds.

## a candidate asymptotic fingerprint

for a representation protocol (r) and register (z), define

[
\phi(z;r)
=========

\Big(
h_{\mathrm{top}},
h,
\alpha,
\beta_h,
\beta_w,
\rho,
c,
d_v,
d_s,
\delta,
\tau,
\operatorname{tw},
e(n),
\eta,
a(n),
q(d),
r_{\mathrm{repair}},
x_{\mathrm{exo}},
m,
v_g
\Big).
]

roughly:

[
h_{\mathrm{top}}
]

is support-growth entropy.

[
h
]

is probabilistic entropy rate.

[
\alpha
]

is the subexponential growth exponent.

[
\beta_h,\beta_w
]

are tree height and width exponents.

[
\rho
]

is branching criticality.

[
c
]

is condensation strength.

[
d_v,d_s
]

are volume-growth and spectral dimensions.

[
\delta
]

is graph hyperbolicity.

[
\tau
]

is the degree-tail exponent or tail family.

[
\operatorname{tw}
]

is semantic treewidth.

[
e(n)
]

is excess-entropy growth.

[
\eta
]

is the long-range correlation exponent.

[
a(n)
]

is ambiguity entropy as a function of length.

[
q(d)
]

is ambiguity survival.

[
r_{\mathrm{repair}}
]

is repair rate and radius.

[
x_{\mathrm{exo}}
]

is exophoric load.

[
m
]

is semantic modularity.

[
v_g
]

is graph-evolution velocity.

this vector would not be “the essence of a language.” it would be an asymptotic phenotype defined relative to a representation and measurement protocol. its value is comparative: it can tell us which distinctions survive scale.

## an empirical research program

first, construct corpora stratified by register: expert-to-expert dialogue, expert-to-novice explanation, scientific papers, administrative documents, ordinary spoken dialogue, fiction, poetry, social media, mathematics, and code.

second, match or control for topic, author, language, token count, and medium. otherwise “scientific versus casual” will be confounded by subject matter, sentence length, transcription conventions, and speaker population.

third, produce several representations for every sample:

a constituency tree.

a dependency tree.

a typed sentence-semantic graph.

a discourse and coreference graph.

a latent-world or event graph.

none should be treated as uniquely correct.

fourth, condition on size. estimate every observable at multiple (n), rather than calculating one corpus-wide average.

fifth, fit finite-size scaling laws:

[
y(n)
====

y_\infty
+
an^{-\omega},
]

or

[
y(n)\sim n^\beta l(n),
]

where (l(n)) is a slowly varying correction. this distinguishes a property that genuinely converges from one that merely appears stable over the observed range.

sixth, compare global and local limits. rescale full trees to study macroscopic shape, but separately inspect fixed-radius neighborhoods around roots, random nodes, discourse hubs, and recently introduced entities.

seventh, estimate the joint distribution of observables. average depth alone may be useless, while the interaction between depth, ambiguity, repair, and semantic treewidth may sharply identify a register.

eighth, train neural models on controlled synthetic languages whose asymptotic properties are known. independently vary branching criticality, ambiguity growth, semantic reentrancy, correlation decay, and graph modularity.

ninth, test whether model internals recover the correct universality class, not only whether output perplexity is low.

tenth, perform interventions. preserve surface statistics while altering causal structure, or preserve semantics while changing grammar presentation. this separates actual invariants from correlations inherited from the dataset.

## the larger conclusion

the deepest possibility is that linguistic universals may not primarily be universal production rules.

they may instead be universal asymptotic constraints.

different human languages, registers, technical domains, and representational systems might use very different local constructions while converging toward a smaller number of large-scale regimes:

critical recursive trees.

spine-dominated narratives.

condensed hub-and-elaboration structures.

small-world semantic graphs.

high-treewidth constraint systems.

low-entropy formal codes.

high-ambiguity interpretive systems.

strongly coupled dialogue processes.

rapidly evolving social graphs.

under this view, a “type of language” is not identified only by its vocabulary or grammar. it is an equivalence class of systems that share the same scaling exponents, limit shapes, information-decay laws, graph topology, ambiguity dynamics, and response to interventions.

the analogy to infinite-width neural networks is therefore valid but limited. infinite-width theory can reveal a kernel or function-space prior, but it does not by itself explain which function finite training selects. likewise, an infinite-language theory could reveal the universality class into which a linguistic system falls, but it would not replace finite syntax, semantics, history, or social context.

its value is that it separates structural law from contingent detail.

instead of asking only:

“what rules does this language use?”

we can ask:

“how does its possibility space grow?”

“how does its hierarchy scale?”

“where does its information live?”

“how long does ambiguity survive?”

“how does meaning diffuse through its graph?”

“which distinctions remain after paraphrase?”

“how much world structure is assumed but unstated?”

“what geometry does expertise create?”

“which observables remain invariant as the language becomes arbitrarily large?”

that is afaict the actual research program hiding inside the original infinite-tree intuition. the tree is only the first slice. the full object is an ensemble of sequences, trees, graphs, worlds, and interacting agents, studied through the transformations and scaling laws that survive in the limit.
